\input{../tex-inputs/latex-leseansicht-vorspann}

\section[Arthur Schnitzler an Gustav Schwarzkopf, 3. 6. 1913]{L04511 Arthur Schnitzler an Gustav Schwarzkopf, 3. 6. 1913}
\nopagebreak\mylabel{L04511v}
\rehead{ }\normalsize\beginnumbering\briefempfaengerindex{Schwarzkopf, Gustav@\textsc{Schwarzkopf, Gustav}!zzzSchnitzler, Arthur@\emph{von Arthur Schnitzler}!1913-06-031@{3. 6. 1913}|(be}
\toendnotes[C]{\smallbreak\pagebreak[2]}
\correspDesc{Versand  durch Arthur Schnitzler am 3. 6. 1913 in Wien
\newline{}Übermittlung  am 3. 6. 1913 in Wien
\newline{}Erhalt  durch Gustav Schwarzkopf im Zeitraum [3. 6. 1913 – 6. 6. 1913?] in Wien}\toendnotes[C]{\smallbreak}
\Standort{DLA, A:Schnitzler, HS.1985.1.1897.}
\physDesc{Bildpostkarte, 294 Zeichen
\newline{}Handschrift: schwarze Tinte, deutsche Kurrent
\newline{}Versand: Stempel: »\nobreak{}\oindex{IX., Alsergrund@\textbf{IX., Alsergrund}|pwk}9/1 Wien 66, 3. VI. 13, 3\nobreak{}«.  }\toendnotes[C]{\smallbreak}
\pstart
           \noindent{}{\pb}\textcolor{gray}{\textbf{Wien, XVIII, Sternwartestr.
                     71\oindex{Wien@\textbf{Wien}!XVIII., Währing@\textbf{XVIII., Währing}!Sternwartestraße 71@\textbf{Sternwartestraße 71}|pw}.}}\pend
           \pstart{}{\pb}Herrn \textsc{Gustav
                     Schwarzkopf}\pend{}\pstart{}\textsc{Baden}\oindex{Baden bei Wien@\textbf{Baden bei Wien}|pw}\pend{}\pstart{}\textsc{bei Wien\oindex{Wien@\textbf{Wien}|pw}};\pend{}\pstart{}\textsc{Pension
                  Quisisana\oindex{Pension Quisisana@\textbf{Pension Quisisana}|pw}},\pend{}\pstart{}Helenenthal\oindex{Helenental@\textbf{Helenental}|pw}.\pend{}{\bigskip}\vspace{1em}
\pstart
           \raggedleft{}3. 6. 913\pend
           \vspace{0.5em}
\pstart
           lieber Guſtav, am \label{K_L04511-1v}\edtext{Freitag Abend}{\lemma{\textnormal{\emph{Freitag Abend}}}\Cendnote{\textnormal{Im \emph{Tagebuch}\pwindex{Schnitzler, Arthur 15. 5. 1862 Wien – 21. 10. 1931 ebd.@\textsc{Schnitzler, Arthur} (15. 5. 1862 Wien – 21. 10. 1931 ebd.)!Tagebuch@\strich\emph{Tagebuch}|pwk} heißt es über die Einladung:
                     »Abend bei uns: Hofrätin Zuckerkandl\pwindex{Zuckerkandl, Berta 13.\,4.\,1864 Wien – 16.\,10.\,1945 Paris@\textsc{Zuckerkandl, Berta} (13.\,4.\,1864 Wien – 16.\,10.\,1945 Paris)|pw} mit Fritz\pwindex{Zuckerkandl, Fritz 30.\,7.\,1895 Wien – 14.\,12.\,1983 Krattigen@\textsc{Zuckerkandl, Fritz} (30.\,7.\,1895 Wien – 14.\,12.\,1983 Krattigen)|pw}, Gound\pwindex{Gund, Robert 18.\,11.\,1865 Neuhausen – 26.\,6.\,1927 Wien@\textsc{Gund, Robert} (18.\,11.\,1865 Neuhausen – 26.\,6.\,1927 Wien)|pw} und Frau (Lauterburg\pwindex{Lauterburg, Elisabeth 18.\,5.\,1885 Langnau – 11.\,11.\,1974 ebd.@\textsc{Lauterburg, Elisabeth} (18.\,5.\,1885 Langnau – 11.\,11.\,1974 ebd.)|pw}), Stephi\pwindex{Bachrach, Stefanie 22.\,5.\,1887 Wien – 16.\,5.\,1917 ebd.@\textsc{Bachrach, Stefanie} (22.\,5.\,1887 Wien – 16.\,5.\,1917 ebd.)|pw}, Mimi\pwindex{Zuckerkandl, Marianne 6.\,8.\,1882 Wien – 1964 Ascona@\textsc{Zuckerkandl, Marianne} (6.\,8.\,1882 Wien – 1964 Ascona)|pw}, Paula Schmidl\pwindex{Schmidl, Paula 13.\,10.\,1874 Wien – 24.\,9.\,1966 Jerusalem@\textsc{Schmidl, Paula} (13.\,10.\,1874 Wien – 24.\,9.\,1966 Jerusalem)|pw}, Else
                           Speidel\pwindex{Speidel-Haeberle, Else 11.\,7.\,1877 Stuttgart – 21.\,7.\,1937 Augustenfeld@\textsc{Speidel-Haeberle, Else} (11.\,7.\,1877 Stuttgart – 21.\,7.\,1937 Augustenfeld)|pw}, Vicky\pwindex{Zuckerkandl, Victor 2.\,7.\,1896 – 24.\,4.\,1965@\textsc{Zuckerkandl, Victor} (2.\,7.\,1896 – 24.\,4.\,1965)|pw}, Liesl\pwindex{Steinrück, Elisabeth 19.\,11.\,1885 – 7.\,4.\,1920 Partenkirchen@\textsc{Steinrück, Elisabeth} (19.\,11.\,1885 – 7.\,4.\,1920 Partenkirchen)|pw}, Albert\pwindex{Steinrück, Albert 20.\,5.\,1872 Wetterburg – 11.\,2.\,1929 Berlin@\textsc{Steinrück, Albert} (20.\,5.\,1872 Wetterburg – 11.\,2.\,1929 Berlin)|pw},
                        Jacobi\pwindex{Jacobi, Bernhard von 27.\,12.\,1880 Hannover – 25.\,10.\,1914 Douai@\textsc{Jacobi, Bernhard von} (27.\,12.\,1880 Hannover – 25.\,10.\,1914 Douai)|pw} und Frau\pwindex{Jacobi, Lucy von 8.\,9.\,1887 Wien – 24.\,4.\,1956 Locarno@\textsc{Jacobi, Lucy von} (8.\,9.\,1887 Wien – 24.\,4.\,1956 Locarno)|pw} (von ›Franziska‹ kommend); ― Erdbeerbowle. Frl. Springer\pwindex{Springer, Gisela 9.\,11.\,1874 Wien – 1942 Vernichtungslager Kulmhof@\textsc{Springer, Gisela} (9.\,11.\,1874 Wien – 1942 Vernichtungslager Kulmhof)|pw}
                        spielte Clavier, Frau Gound\pwindex{Lauterburg, Elisabeth 18.\,5.\,1885 Langnau – 11.\,11.\,1974 ebd.@\textsc{Lauterburg, Elisabeth} (18.\,5.\,1885 Langnau – 11.\,11.\,1974 ebd.)|pw}, Olga\pwindex{Schnitzler, Olga 17.\,1.\,1882 Wien – 13.\,1.\,1970 Lugano@\textsc{Schnitzler, Olga} (17.\,1.\,1882 Wien – 13.\,1.\,1970 Lugano)|pw} sangen; Albert\pwindex{Steinrück, Albert 20.\,5.\,1872 Wetterburg – 11.\,2.\,1929 Berlin@\textsc{Steinrück, Albert} (20.\,5.\,1872 Wetterburg – 11.\,2.\,1929 Berlin)|pw} sang komisches. Es war anfangs
                     leidlich, dauerte aber zu lang. (1 Uhr.)«, siehe A. S.: \emph{Tagebuch}, 6. 6. 1913.
               }}}\label{K_L04511-1} haben wir etliche Bekannte\pwindex{Zuckerkandl, Berta 13.\,4.\,1864 Wien – 16.\,10.\,1945 Paris@\textsc{Zuckerkandl, Berta} (13.\,4.\,1864 Wien – 16.\,10.\,1945 Paris)|pwv}\pwindex{Zuckerkandl, Fritz 30.\,7.\,1895 Wien – 14.\,12.\,1983 Krattigen@\textsc{Zuckerkandl, Fritz} (30.\,7.\,1895 Wien – 14.\,12.\,1983 Krattigen)|pwv}\pwindex{Gund, Robert 18.\,11.\,1865 Neuhausen – 26.\,6.\,1927 Wien@\textsc{Gund, Robert} (18.\,11.\,1865 Neuhausen – 26.\,6.\,1927 Wien)|pwv}\pwindex{Lauterburg, Elisabeth 18.\,5.\,1885 Langnau – 11.\,11.\,1974 ebd.@\textsc{Lauterburg, Elisabeth} (18.\,5.\,1885 Langnau – 11.\,11.\,1974 ebd.)|pwv}\pwindex{Bachrach, Stefanie 22.\,5.\,1887 Wien – 16.\,5.\,1917 ebd.@\textsc{Bachrach, Stefanie} (22.\,5.\,1887 Wien – 16.\,5.\,1917 ebd.)|pwv}\pwindex{Zuckerkandl, Marianne 6.\,8.\,1882 Wien – 1964 Ascona@\textsc{Zuckerkandl, Marianne} (6.\,8.\,1882 Wien – 1964 Ascona)|pwv}\pwindex{Schmidl, Paula 13.\,10.\,1874 Wien – 24.\,9.\,1966 Jerusalem@\textsc{Schmidl, Paula} (13.\,10.\,1874 Wien – 24.\,9.\,1966 Jerusalem)|pwv}\pwindex{Speidel-Haeberle, Else 11.\,7.\,1877 Stuttgart – 21.\,7.\,1937 Augustenfeld@\textsc{Speidel-Haeberle, Else} (11.\,7.\,1877 Stuttgart – 21.\,7.\,1937 Augustenfeld)|pwv}\pwindex{Zuckerkandl, Victor 2.\,7.\,1896 – 24.\,4.\,1965@\textsc{Zuckerkandl, Victor} (2.\,7.\,1896 – 24.\,4.\,1965)|pwv}\pwindex{Steinrück, Albert 20.\,5.\,1872 Wetterburg – 11.\,2.\,1929 Berlin@\textsc{Steinrück, Albert} (20.\,5.\,1872 Wetterburg – 11.\,2.\,1929 Berlin)|pwv}\pwindex{Jacobi, Bernhard von 27.\,12.\,1880 Hannover – 25.\,10.\,1914 Douai@\textsc{Jacobi, Bernhard von} (27.\,12.\,1880 Hannover – 25.\,10.\,1914 Douai)|pwv}\pwindex{Jacobi, Lucy von 8.\,9.\,1887 Wien – 24.\,4.\,1956 Locarno@\textsc{Jacobi, Lucy von} (8.\,9.\,1887 Wien – 24.\,4.\,1956 Locarno)|pwv} bei uns und \textsc{Liesl\pwindex{Steinrück, Elisabeth 19.\,11.\,1885 – 7.\,4.\,1920 Partenkirchen@\textsc{Steinrück, Elisabeth} (19.\,11.\,1885 – 7.\,4.\,1920 Partenkirchen)|pw}};  es wäre{ }ſehr{ }ſchön von Ihnen, we{\geminationn} Sie, ſtatt des \label{K_L04511-2v}\edtext{projekt. Nachmittags}{\lemma{\textnormal{\emph{projekt. Nachmittags}}}\Cendnote{\textnormal{Schwarzkopf hatte
                  seinen Besuch für den Nachmittag des 6. 6. 1913
                  angekündigt, siehe XXXX Auszeichnungsfehler: Dokument L04257 nicht gefunden.}}}\label{K_L04511-2}, uns am Abend die Freude machten{[}.{]}
               Natürlich früh. (7 – ½ 8). –\pend
           \pstart Herzlichſt Ihr \spacefill\mbox{Arthur}\pend{}\selectlanguage{ngerman}\endnumbering\briefempfaengerindex{Schwarzkopf, Gustav@\textsc{Schwarzkopf, Gustav}!zzzSchnitzler, Arthur@\emph{von Arthur Schnitzler}!1913-06-031@{3. 6. 1913}|)be}\mylabel{L04511h}
\begin{anhang}
\end{anhang}\newcommand{\dateiname}{L04511}\newcommand{\titel}{Arthur Schnitzler an Gustav Schwarzkopf, 3. 6. 1913}\newcommand{\editorInnen}{Herausgegeben von Jahnke, SelmaMüller, Martin Anton}\input{../tex-inputs/latex-leseansicht-abspann}
